%%
%% Patterns (Cell Press) Manuscript Shell (Option C Fallback)
%% Venue: Patterns (Cell Press)
%% Submission Mode: Rolling submission (contingent on post-rejection activation from Nature MI)
%%
\documentclass[11pt,letterpaper]{article}

\usepackage[utf8]{inputenc}
\usepackage{amsmath,amssymb,amsfonts}
\usepackage{graphicx}
\usepackage{booktabs}
\usepackage{lineno}
\usepackage{cite}
\usepackage{hyperref}
\usepackage[margin=1in]{geometry}

\linenumbers

\title{\textbf{Fragility of Download-Weighted Safety Metrics on Public Model Registries: An Empirical Census and Heavy-Tail Analysis}}

\author{
  Petar Radanliev$^{1,2,3,*}$\\[1ex]
  \small $^{1}$School of Computing and Communications, STEM, The Open University, Milton Keynes, UK\\
  \small $^{2}$The Alan Turing Institute, British Library, London, UK\\
  \small $^{3}$Department of Computer Science, University of Oxford, Oxford, UK\\
  \small $^*$Correspondence: \texttt{petar.radanliev@open.ac.uk}
}

\date{\today}

\begin{document}

\maketitle

\section*{The Bigger Picture}
Public model hubs function as modern open-source software registries for artificial intelligence. However, documentation debt—missing behavioral boundaries, safety disclosures, and licensing—poses severe supply-chain risks. While stakeholders often rely on download counts as an exposure weighting proxy, our empirical census demonstrates that download-weighted metrics follow an infinite-variance Pareto distribution, making raw metrics brittle and dominated by corporate outliers. Robust data stewardship requires tail-trimmed indicators and supply-chain provenance tracking.

\section*{Highlights}
\begin{itemize}
  \item Census of 809,703 model cards and 48,247 lineage edges across 2.47 billion downloads.
  \item Popularity dividend is falsified on the registered estimator ($-2.9$ percentage points).
  \item Extreme Value Theory establishes an infinite-variance Pareto regime ($\alpha = 0.98 \le 2$).
  \item Single top model shifts global download-weighted safety rate by 7.20 percentage points.
  \item Model merges display an 83.46\% license omission rate across supply chains.
\end{itemize}

\begin{abstract}
Public model repositories constitute critical digital infrastructure for distributing open-weight artificial intelligence models. While model cards serve as primary instruments for algorithmic transparency, whether documentation practices scale with model usage remains an open empirical question. Here, we present a comprehensive census of safety documentation across 809,703 model cards and 48,247 supply-chain lineage edges on the Hugging Face Hub, encompassing 2.47 billion downloads. Evaluating a pre-registered popularity hypothesis, we find the popularity gap is falsified on the registered binned estimator ($-2.9$ percentage points). On the unregistered row-level vector, the download-weighted safety rate (38.40\%) exceeds unweighted rates (27.19\%) by 11.21 points, but the 95\% bootstrap interval ([$ -1.67$, $+24.68$] points) spans both boundaries and is indeterminate. Extreme Value Theory indicates repository downloads follow an infinite-variance Pareto distribution (Hill tail index $\alpha = 0.98 \le 2$), causing weighted metrics to be hypersensitive to individual corporate releases: omitting the top model alone swings the weighted safety rate by 7.20 percentage points. Furthermore, derivative lineages exhibit severe documentation decay, with an 83.46\% license omission rate on merges. These findings reveal structural vulnerabilities in naive repository metrics, establishing the necessity of tail-trimmed estimators and provenance tracking for data stewardship in open AI ecosystems.
\end{abstract}

\section*{Introduction}
Machine learning model cards formalise documentation across public repositories. In this study, we investigate whether download-weighted safety disclosure reflects an organic market dividend or statistical fragility.

\section*{Results}
\subsection*{Census Baseline and Falsification of the Popularity Dividend}
Across 809,703 model cards, unweighted safety disclosure ($L_2+$) is 27.19\%. Evaluating the pre-registered binned popularity model yields a gap of $-2.9$ points (falsified). Unregistered row-level weighting produces an apparent $+11.21$ point gap, with a 95\% bootstrap interval of [$-1.67$, $+24.68$] points.

\subsection*{Pareto Heavy-Tail Dynamics and Outlier Leverage}
Hill estimation confirms an infinite-variance regime ($\alpha = 0.98$). Influence function calculations show that removing the single most downloaded model alters the global rate by 7.20 percentage points.

\subsection*{Lineage Decay and Licensing Friction}
Across 48,247 derivative edges, child models exhibit significant documentation decay. Merged models omit upstream license declarations in 83.46\% of cases.

\section*{Discussion}
Data stewardship in foundation model ecosystems cannot rely on raw download metrics. We discuss implications for AI governance, dataset archiving, and transparent data repositories.

\section*{Experimental Procedures}
\subsection*{Corpus Extraction and Adjudication}
Repository snapshots were processed via hierarchical YAML and Markdown parsing. Construct validity was established through human annotation across stratified subsets ($\kappa = 0.81$).

\subsection*{Statistical and EVT Analysis}
Heavy-tail indices were estimated via upper-order statistic Hill regressions. Confidence intervals were computed using 1,000 bootstrap replicates over the empirical Lorenz curve.

\bibliographystyle{cell}
\bibliography{manuscript}

\end{document}
